\newcommand\lword[1]{\leavevmode\nobreak\hskip0pt plus\linewidth\penalty50\hskip0pt plus-\linewidth\nobreak #1}
\newcommand\pill[2][]{\lword{\tikz[overlay]\node[fill=#1,inner sep=2pt, anchor=text, rectangle, rounded corners=1mm]{#2};\phantom{#2}}}

\newcommand{\kpill}[1]{\pill[mila-khaki]{#1}}


\newcommand\usfcomponent[4][]{
    \begin{tcolorbox}[
        colframe=#1,
        colback=#1,
        size=fbox,
        enhanced,
        coltitle=black, 
        detach title,
        left=5mm,
        overlay={
            \node[xshift=2mm] at (frame.west) {\pill[#2]{\tcbtitle}};
        },
        title=#3
        ]
        {#4}        
    \end{tcolorbox}
}

\newcommand{\tbunifiedscratch}{
    \begin{wrapfigure}[25]{r}{0.35\textwidth}
        \raisebox{0pt}[\dimexpr\height-1\baselineskip\relax]{
            \centering
            \scriptsize
            \begin{minipage}{\linewidth}
                \tcbset{parskip}
                \begin{framed}
                    \textcolor{gray}{Input ($\vx$)}\\
                    \texttt{Compute 5 3 7 2 6 + 1 9 1 7 =}
                    \vspace{0.4em} \\
                    \textcolor{gray}{Output ($\vs;\vy$)}\\
                    {
                    \ttfamily
                    <scratch>
                    \vspace{0.5em}
                    \usfcomponent[mila-purple!10!white]{mila-purple!30!white}{$\mathcal{I}$}{
                        For digits 6 and 7,
                    }
                    \vspace{0.5mm}
                    \usfcomponent[mila-khaki!50!white]{mila-khaki}{$\mathcal{C}$}{
                        We have ( 6 + 7 + carry ) \% 10 =
                        13 \% 10 = 13 \% 10
                    }
                    \vspace{0.5mm}
                    \usfcomponent[plot-blue!25!white]{plot-blue!60!white}{$\mathcal{O}$}{
                        Which is equal to 3 .
                    }
                    \vspace{0.5mm}
                    \usfcomponent[plot-red!25!white]{plot-red!60!white}{$\mathcal{V}$}{
                        We update carry to 13 // 10 = 1.
                    }
                    \vspace{0.5mm}
                    \usfcomponent[plot-green!25!white]{plot-green!60!white}{$\mathcal{R}$}{
                        So, the remaining input is 5 3 7 2 + 1 9 1
                    }
                    \vspace{0.5em}
                    \dots
                    \vspace{0.5em}
                    \\
                    </scratch> \\
                    The answer is 5 5 6 4 3.

                    }
                \end{framed}
                \caption{Example of an addition task depicted with its first scratchpad step.
                    Each step consists of five components:
                    \ \pill[mila-purple!30!white]{\phantom{A}} Step Input $\mathcal{I}$,
                    \ \pill[mila-khaki]{\phantom{A}} Step Computation $\mathcal{C}$,
                    \ \pill[plot-blue!60!white]{\phantom{A}} Step Output $\mathcal{O}$,
                    \ \pill[plot-red!60!white]{\phantom{A}} Intermediate Variable Updates  $\mathcal{V}$,
                    and \ \pill[plot-green!60!white]{\phantom{A}} Remaining Input $\mathcal{R}$.
                }
                \label{fig:unified_scratchpad_format}
            \end{minipage}
        }
    \end{wrapfigure}
}